% !TeX root = ../MAgent_Theory_2.0_companion.tex
\chapter{Two-Stage Selection on a Wishart Family}
\label{cf:formation}

This chapter concerns the reciprocal law potential, the direct interaction potential between agent laws whose Fisher gradient flow moves each source toward the receivers that attend to it. It is neither the global free energy nor the per-agent free energy of the dynamics chapter, and nothing here is transferred to those objectives. Four distinct agents are grouped into two parents and then into a parent of those parents, while each parent keeps the original smooth KL/Fisher dynamics, a declared response interface, and the collective connection of its children.

The example is best read as exact inheritance along a prepared ordering. The rule that groups the agents reads their belief shapes, and every belief shape relaxes as $\dot\alpha_i=-2(\alpha_i-\alpha_0)$, with no trace of the coupling, of the precisions, of the model sector, or of any other agent. The grouping is therefore the centroid dendrogram of the initial shapes, read on a clock whose origin the threshold sets, and the same grouping occurs, at the same times, at zero coupling. The formation standard of the inheritance chapter compares the flow with two null flows. In the renormalized null, which removes the peer edges together with their labels, the grouping is reproduced exactly. In the label-preserving null, which keeps the labels and priors, each agent's self row loses weight to the empty peer entries by an amount that depends on how far the agent is from its reference, the agents relax on different clocks, and integrations show a different grouping on an open set of preparations, already at zero coupling. That difference has nothing to do with the interaction, the two verdicts disagree, and the standard counts interaction necessity as failed. The second stage also forms from every counterfactual handoff, so the construction is a two-stage selection with exact inheritance, not a recursive formation. The architecture assumes common matrix precisions, a constant model-independent evaluator, a particular moment-dependent interaction, fresh sum responses, and a fixed membership threshold. All block-sum descriptions already exist mathematically at the initial time, and they are exact for every partition; the threshold selects when one becomes operative, and the dynamics does not first create the redundancy that permits it. The initial class is open within this statistical architecture.

\section{Random matrices as represented contents}
\label{cf:laws}

Let both represented-content spaces be $\mathsf K=\mathsf M=\operatorname{SPD}(2)$, the cone of real symmetric positive-definite two-by-two matrices. A sample $k$ or $m$ is a matrix content. It is distinct from a contextual coordinate and need not be the covariance of another random variable. Use the Wishart density
\begin{equation}
 \begin{aligned}
 W_2(\nu,\mathsf P)(dk)
 &=\frac{\det(\mathsf P)^\nu}{\Gamma_2(\nu)}
 \det(k)^{\nu-3/2}e^{-\operatorname{tr}(\mathsf Pk)}dk,\\
 \Gamma_2(\nu)&=\sqrt\pi\Gamma(\nu)\Gamma(\nu-\tfrac12),
 \end{aligned}
 \label{cf:wishart}
\end{equation}
where $\nu>1/2$, $\mathsf P\succ0$, and $dk$ is Lebesgue measure on the three independent entries. The shape $\nu$ is real; a Gaussian sampling representation is not assumed. $\Gamma_2$ is the multivariate gamma function \citep{Muirhead1982}, and $\nu>1/2$ is exactly the range of real shapes for which this density is a probability law \citep{Gindikin1975}; \citet{LetacMassam2017} discuss the harder noncentral case.

Normalization can be checked by writing $k_{11}=u$, $k_{12}=\sqrt u v$, and $k_{22}=w+v^2$, with $u,w>0$. At $\mathsf P=I_2$, the determinant is $uw$, the trace is $u+w+v^2$, and the Jacobian is $\sqrt u$. Two gamma integrals and one scalar Gaussian integral give $\Gamma_2(\nu)$. Congruence by $\mathsf P^{-1/2}$ gives the general normalizer. The same calculation yields
\begin{equation}
 \E e^{-\operatorname{tr}(Tk)}
 =\left(\frac{\det\mathsf P}{\det(\mathsf P+T)}\right)^\nu,
 \qquad
 \E k=\nu\mathsf P^{-1},
 \qquad \mathsf P+T\succ0.
 \label{cf:laplace}
\end{equation}
Multiplying these transforms proves that independent matrices with common precision add their shapes. Transform uniqueness applies because the transforms are finite near $T=0$. Thus $k_i\sim W_2(\alpha_i,\mathsf P)$ independently implies $\sum_{i\in I}k_i\sim W_2(\sum_{i\in I}\alpha_i,\mathsf P)$. This is an exact convolution identity for laws of matrices.

Take four agents with $q_i=W_2(\alpha_i,\mathsf P)$, $s_i=W_2(\eta_i,\mathsf Q)$, and product recognition $\zeta_i=q_i\otimes s_i$. Shapes exceed $1/2$. They may differ across agents and preparations, but are constant over the context base. They will evolve in inference time. The common precisions are preparation assumptions.

\section{An interaction that moves every precision and no shape}
\label{cf:flow}

Fix $\alpha_0,\eta_0>1/2$ and $\kappa>0$. The proper local reference is the product of $W_2(\alpha_0,I_2)$ and $W_2(\eta_0,I_2)$; its normalized evaluator is the constant kernel $G_i(m,dk)=W_2(\alpha_0,I_2)(dk)$. The constancy matters. A belief evaluator whose precision depended on the model sample $m$ would add the average of that precision over $s_i$ to the force on $q_i$; since the model laws differ across agents, their precisions would separate and the sum identity above would no longer produce Wishart parents. With the constant evaluator the two sectors are uncoupled. Let $N=4$ remain the primitive population count. Define
\begin{equation}
 \begin{aligned}
 M_q&=\sum_i\E_{q_i}k,
 \qquad \Lambda_q=I_2+\frac\kappa N M_q,\\
 \mathcal K^q_{ij}[\boldsymbol q](k',dk)
 &=W_2(\alpha_0,\Lambda_q)(dk).
 \end{aligned}
 \label{cf:channel}
\end{equation}
The model sector exchanges $(q,\alpha_0,k)$ for $(s,\eta_0,m)$. Every source, including self, is admitted with prior $\pi_{ij}=1/N$; the self channel is deliberately not the identity, since an identity self channel would have zero cost and would tilt every row away from its prior. These normalized channels ignore the individual input sample $k'$, but read the population's current moment. Each receiver's costs are identical across source labels, so exact Gibbs reduction at any positive source temperature gives the original prior row. The reduced belief potential is
\begin{equation}
 \mathcal E_q=\sum_i\left[
 \KL(q_i\Vert W_2(\alpha_0,I_2))
 +\KL(q_i\Vert W_2(\alpha_0,\Lambda_q))\right].
\end{equation}

Changing one law changes every target in \cref{cf:channel}. The moment-dependent peer terms are $\operatorname{tr}(\Lambda_qM_q)-N\alpha_0\log\det\Lambda_q$. Their full differential is $\operatorname{tr}(C_qdM_q)$, where
\begin{equation}
 C_q=\Lambda_q+\frac\kappa N
 (M_q-N\alpha_0\Lambda_q^{-1})
 =I_2+\frac{2\kappa}{N}M_q-\kappa\alpha_0\Lambda_q^{-1}.
 \label{cf:channel-force}
\end{equation}
The second term in the first expression comes from differentiating the channel. Freezing the target would remove it. Up to a constant, the density force and its ordinary unit-Fisher flow are
\begin{equation}
 \begin{aligned}
 h_i(k)&=2\log q_i(k)-2(\alpha_0-\tfrac32)\log\det k
       +\operatorname{tr}[(I_2+C_q)k],\\
 \dot q_i&=-q_i(h_i-\E_{q_i}h_i),\\
 \dot\alpha_i&=-2(\alpha_i-\alpha_0),
 \qquad \dot{\mathsf P}=-2\mathsf P+I_2+C_q.
 \end{aligned}
 \label{cf:full-flow}
\end{equation}
The last two equations follow by substituting \cref{cf:wishart} into the density equation and matching the coefficients of $\log\det k$ and $k$. They establish tangency of the unrestricted density field. No parameter projection or mobility factor is inserted. The model equations use $(\eta_i,\eta_0,\mathsf Q,C_s)$.

The matching shows where the interaction acts. Substituting $\log q_i=(\alpha_i-\tfrac32)\log\det k-\operatorname{tr}(\mathsf Pk)+\mathrm{const}$ turns the force into $2(\alpha_i-\alpha_0)\log\det k$ plus a term linear in $k$. The coupling $\kappa$ and the other agents live only in the linear term, so they move the precision and never the shape. The precisions are genuinely coupled: a moment perturbation $H$ gives $DC_q[H]=(2\kappa/N)H+(\kappa^2\alpha_0/N)\Lambda_q^{-1}H\Lambda_q^{-1}$, whose trace pairing with $H\ne0$ is positive. The product evaluator supplies no coupling between sectors.

A general picture lies behind this. When every receiver sees the same evaluator and the same channel target, and the target reads the laws only through an additive moment, the force on each agent is $2\log q_i$ plus one function common to all agents. Every log-likelihood ratio between two agents then decays as $e^{-2t}$ whatever the interaction, and in an exponential family every difference of natural parameters does too. The precision difference $\mathsf P_i(t)-\mathsf P_j(t)=e^{-2t}(\mathsf P_i(0)-\mathsf P_j(0))$ and the shape law are two instances. A grouping rule that compares agents only through these log-likelihood ratios can only see a shrinking copy of the starting configuration. A rule that reads absolute quantities of the potential, such as each agent's cost against the common target, sees the coupling through one common scalar and can depend on it; whether the interaction matters for grouping is therefore a statement relative to the declared rule. On finite tables the same identity is a case of the common-clock proposition of the inheritance chapter, which also states what happens without the peer costs.

These equations remain regular for all finite times. Put $\bar\alpha=\sum_i\alpha_i/N$. Then $M_q=N\bar\alpha\mathsf P^{-1}$ and each precision eigenvalue satisfies
\begin{equation}
 \dot p=f(p,\bar\alpha(t))
 =-2p+2+\frac{2\kappa\bar\alpha(t)}p
 -\frac{\kappa\alpha_0p}{p+\kappa\bar\alpha(t)}.
 \label{cf:eigenvalue}
\end{equation}
The mean shape relaxes to $\alpha_0$ and stays above $1/2$. The right side is positive for sufficiently small $p$ and negative for sufficiently large $p$, uniformly along this shape trajectory. Each eigenvalue therefore remains in a compact positive interval. Shapes also remain in compact intervals above $1/2$, ensuring finite KL values and square-integrable scores. Common precision is preserved, and starting with unequal precisions they never become equal at a finite time.

\section{The same clock after joining}
\label{cf:parent}

For a current block $I$, retain $n_I=|I|$, $A_I=\sum_{i\in I}\alpha_i$, $B_I=\sum_{i\in I}\eta_i$, and the common precisions. Fresh independent addition gives $q_I=W_2(A_I,\mathsf P)$ and $s_I=W_2(B_I,\mathsf Q)$. The parent references and evaluator are the convolutions of the children's references and evaluator outputs; for the evaluator this is a kernel from the parent's model sum only because the child evaluators ignore the model sample. Their shapes are $n_I\alpha_0,n_I\eta_0$; the parent peer target has belief shape $n_I\alpha_0$ and precision $\Lambda_q$. Its source prior is $\pi^c_{IJ}=n_J/N$, the sum of the original source probabilities over sender block $J$.

The total moment is still $M_q=\sum_I\E_{q_I}k_I$, and $N=\sum_In_I$ does not become the number of current parents. The inherited original-form potential is
\begin{equation}
 \mathcal E_q^c=\sum_I\left[
 \KL(q_I\Vert W_2(n_I\alpha_0,I_2))
 +\KL(q_I\Vert W_2(n_I\alpha_0,\Lambda_q))\right].
\end{equation}
Its target normalizers still contribute $-N\alpha_0\log\det\Lambda_q$. The complete derivative is consequently the same $C_q$, giving the unit-Fisher equations
\begin{equation}
 \dot A_I=-2(A_I-n_I\alpha_0),
 \qquad \dot B_I=-2(B_I-n_I\eta_0),
 \qquad \dot{\mathsf P}=-2\mathsf P+I_2+C_q,
 \label{cf:parent-flow}
\end{equation}
and the exchanged equation for $\mathsf Q$. These are exactly the summed fine shape equations and original precision flow, in the same time. Fine and parent energies need not be equal. Joining adds $(n,A,B)$ and keeps the common precisions, so joining is associative and commutes with evolution. Independent addition also commutes with the evaluator and channel outputs. All of this holds for every partition, so inheritance alone does not single out the blocks that the rule below picks. Even the final single parent keeps its inherited nonidentity self channel and second KL term; deleting it would change the flow.

\section{Two selected parents and an uncoupled lost allocation}
\label{cf:selection}

Fix a threshold $\delta>0$ and start with singletons whose belief shapes are pairwise farther than $\delta$ apart. A block's score is $b_I=A_I/n_I-\alpha_0$. It obeys $\dot b_I=-2b_I$; joining replaces two scores by their count-weighted average. Join current blocks when $|b_I-b_J|$ first reaches $\delta$ from above. Simultaneous threshold edges join their connected components. This is a declared, relabeling-equivariant rule applied to current states. Two shapes that start closer than $\delta$ would never be joined, because their distance never reaches $\delta$ from above; the separation assumption is part of the rule.

Because every score is the initial centroid times $e^{-2t}$, a pair of blocks reaches the threshold at $\tfrac12\log(D/\delta)$, where $D$ is the distance between their initial centroids. The first join is the closest pair. After it, the merged centroid lies between its two parts, and every other block lies outside that segment, so its new distance exceeds the merge distance. Heights therefore increase, nothing cascades, and the rule reproduces count-weighted centroid clustering of the initial shapes, the UPGMC method of numerical taxonomy \citep{SokalMichener1958,LanceWilliams1967}, with join $k$ at time $\tfrac12\log(D_k/\delta)$. It is not single linkage: for shapes $0,1,2.05,3.5$ single linkage attaches the third point to the first pair, while the rule forms the two pairs first. Neither $\kappa$, nor the precisions, nor the model shapes enter, and changing $\delta$ only shifts every time by $-\tfrac12\log\delta$. At zero coupling the same events occur at the same times, and so they do when the peer edges are removed together with their labels and each self channel reads the other agents' initial moment. If instead the labels and priors are kept, the self rows are reweighted, the clocks become agent dependent, and the events change, also at zero coupling.

Sort the initial belief shapes as $u_1<u_2<u_3<u_4$, and let $d_1,d_2,d_3$ be successive gaps. The lower pair and the upper pair are joined first, followed by their union, exactly when
\begin{equation}
 \delta<\min(d_1,d_3)<d_2,
 \qquad \max(d_1,d_3)<d_2+\tfrac12\min(d_1,d_3),
 \label{cf:gaps}
\end{equation}
up to ties. The first two times and the union time are
\begin{equation}
 \begin{aligned}
 t_L&=\tfrac12\log(d_1/\delta),
 \qquad t_R=\tfrac12\log(d_3/\delta),\\
 t_2&=\tfrac12\log\frac{d_2+(d_1+d_3)/2}{\delta}.
 \end{aligned}
 \label{cf:times}
\end{equation}
To see why no other join intervenes, suppose the lower pair joins first. Its closest remaining singleton has distance $e^{-2t}(d_2+d_1/2)$, larger than the upper pair's distance $e^{-2t}d_3$ by the second condition. Thus the upper pair joins next. The reversed ordering has the same argument. The resulting parent separation is $e^{-2t}[d_2+(d_1+d_3)/2]$, the centroid distance rather than the nearest-neighbor gap $d_2$, so both first parents coexist for a positive interval. Each selected child-block distance remains below threshold afterward. All primitive shape differences remain nonzero at every finite time, but every law converges to one common limit: the parents are distinct at finite times and the population approaches consensus asymptotically.

For explicit values choose $\alpha_0=\eta_0=\kappa=1$, $\delta=0.1$, and initial belief shapes $(1,1.2,1.8,2.1)$. The gaps are $(0.2,0.6,0.3)$ and the times are $\tfrac12\log2$, $\tfrac12\log3$, and $\tfrac12\log(17/2)$; the first parents coexist for $\tfrac12\log(17/6)\approx0.52$. Model shapes can vary independently near $(0.8,1.1,1.4,1.7)$. A small open neighborhood and all its relabelings preserve the strict inequalities. No binary tree was supplied; all additive block descriptions were available before any threshold crossing.

A formation time can reveal a belief-shape gap: at a join the two child scores differ by exactly $\delta$, so the parent sum, the counts, and the time recover the children's belief shapes up to exchange. To exhibit information that remains lost even with all times recorded, fix the belief shapes and change two model shapes inside a joining pair by $(\eta_i,\eta_j)\mapsto(\eta_i+\epsilon,\eta_j-\epsilon)$. The parent model shape, population moments, subsequent precision flow, and all selection times stay unchanged. The individual public model laws change continuously, with a nonzero Fisher tangent because their shape scores contain the nonconstant statistic $\log\det m$. Distinct fixed belief labels exclude a mere permutation. At the second join, transfer model shape from one first-parent sum to the other while fixing the total. The final parent loses a further continuous distinction between its children. The lost allocation is, however, a spectator even before the join: each model shape relaxes on its own, the model precision sees only the total, the belief sector never reads the model laws, and the rule reads only belief shapes. The loss shows that the parent description is many-to-one, not that it discards something any agent responded to.

\section{What a parent can be asked}
\label{cf:responses}

At each query, draw fresh independent samples $k_i\sim q_i$, $m_i\sim s_i$ across agents and sectors. For every current block report $k_I=\sum_{i\in I}k_i$, $m_I=\sum_{i\in I}m_i$. At partition $\Pi$, the whole response vector has the conditional law
\begin{equation}
 K_\Pi=\bigotimes_{I\in\Pi}
 \left[W_2(A_I,\mathsf P)\otimes W_2(B_I,\mathsf Q)\right].
 \label{cf:response-kernel}
\end{equation}
Allow any normalized readout kernel on this vector. This includes overlapping unions of whole current blocks using the same batch of samples, so their correlations are retained. A controller may choose future readouts and finite query times from its carried record history and current retained state.

Condition on that history. Fine and parent waiting-time evolution agree by \cref{cf:parent-flow}, joins give the same retained state, and \cref{cf:response-kernel} gives the same next-record kernel. Induction therefore gives the same finite joint history of records, actions, times, and retained states for each initial preparation. The argument is short for a reason: the interface sees only block sums, the retained state is a deterministic function of preparation and time, and records never feed back into the laws, so adaptivity only chooses when and what to read. Any expected cost or record entropy of an adaptive policy is therefore the same in both descriptions. Earlier records are carried forward; different preparations need not have identical record laws from time zero. Recording a response does not reset any belief. A query cutting through a retained block can reveal an unretained intersection shape. Reusing a hidden child's samples, applying child-specific Bayesian updates, or letting records update the laws requires another interface and another closure proof.

\section{Prepared curvature that survives both joins and then fades}
\label{cf:geometry}

Use the contextual chart $(a,\chi)\in(1/5,2/5)\times(-\chi_0,\chi_0)$, with $\chi_0>0$, and one $SO(2)$ acting by matrix conjugation on belief and model samples together. Let $R_\chi$ be rotation through angle $\chi$, and prepare
\begin{equation}
 \mathsf P(0;a,\chi)=R_\chi
 \operatorname{diag}(e^a,e^{-a})R_\chi^T,
 \qquad
 \mathsf Q(0)=\operatorname{diag}(e^{3/10},e^{-3/10}).
 \label{cf:section}
\end{equation}
These coordinates describe how the laws depend on context, whereas $k,m$ remain represented matrix samples. Context-constant shapes make all event times independent of context. The flow acts separately at each context, and nothing in it reads a connection: the connection below is computed from the evolved sections as a diagnostic. At fixed shape, a precision variation $U$ has score and Fisher metric
\begin{equation}
 \begin{aligned}
 \ell_U(k)&=\nu\operatorname{tr}(\mathsf P^{-1}U)-\operatorname{tr}(Uk),\\
 g^F(U,V)&=\nu\operatorname{tr}(\mathsf P^{-1}U\mathsf P^{-1}V).
 \end{aligned}
 \label{cf:score}
\end{equation}
The joint score for common contextual changes depends on the children only through their matrix sums. The same holds for common rotation scores. Addition therefore retains every contextual and orbit direction used here, with its exact Fisher inner products. Fine and parent infinitesimal stabilizers also agree. With unit child weights, the collective minimizing connection and residual contextual tensor are identical before and after a join.

Write $p_+>p_-$ for the belief precision eigenvalues, and $\tilde p_+>\tilde p_-$ for the model precision eigenvalues. The derivative $f_p$ in \cref{cf:eigenvalue} is negative, and the scalar solution's derivative with respect to its initial value is $\exp(\int_0^t f_pdu)>0$ at finite time. Thus $\partial_ap_+>0$, $\partial_ap_-<0$, and neither eigenvalue pair meets in finite time. For a block $I$, define its two rotational inertias and radial coefficient by
\begin{equation}
 \begin{aligned}
 \mathcal I_1&=2A_I\frac{(p_+-p_-)^2}{p_+p_-},
 &\mathcal I_2&=2B_I\frac{(\tilde p_+-\tilde p_-)^2}{\tilde p_+\tilde p_-},\\
 g_{aa}&=A_I\left[
 (\partial_ap_+/p_+)^2+(\partial_ap_-/p_-)^2\right].
 \end{aligned}
\end{equation}
All three are positive, $\partial_a\mathcal I_1>0$, and $\mathcal I_2$ is independent of $a$. An angular connection correction $\xi$ has Fisher cost $\mathcal I_1(d\chi+\xi)^2+\mathcal I_2\xi^2$. Minimization, with the convention $D^{\mathcal A}p=dp+J_p\mathcal A$, gives
\begin{equation}
 \begin{aligned}
 \mathcal A^F&=-\frac{\mathcal I_1}{\mathcal I_1+\mathcal I_2}d\chi,
 &h^{\min}&=g_{aa}da^2+\frac{\mathcal I_1\mathcal I_2}{\mathcal I_1+\mathcal I_2}d\chi^2,\\
 \mathcal F&=-\frac{(\partial_a\mathcal I_1)\mathcal I_2}{(\mathcal I_1+\mathcal I_2)^2}da\wedge d\chi
 \ne0.
 \end{aligned}
 \label{cf:curvature}
\end{equation}
The scalar connection form represents the coefficient of the standard $SO(2)$ generator. Its curvature is its exterior derivative. The tensor has rank two at every finite time. Applying the same formulas to each block's fine children or retained parent gives equality through both joins. A passive frame rotation by $\vartheta$ sends $\mathcal A^F$ to $\mathcal A^F-d\vartheta$, with the compensating change in the section derivative. The fixed addition map contributes no moving-channel term to this geometric comparison. Physically, this is the gauge field of a two-part rotor whose belief part has an inertia that varies along the context while the model part does not; it is the mechanical connection of rotating bodies in a statistical setting \citep{Smale1970}.

Three qualifications keep this picture honest. First, the curvature exists because one rotation is declared to act on beliefs and models together. Nothing in the potential ties them; with separate rotations for the two sectors, the minimizing connection is $-d\chi$ on the belief factor and zero on the model factor, which is flat, and the tensor has rank one. Second, the contextual dependence was put into the initial section, and none of the formulas uses $\kappa>0$: at zero coupling the same curvature appears. Third, the flow erases it. All contexts relax to the same isotropic precision $p_*I_2$, where $p_*$ is the positive zero of $f(\cdot,\alpha_0)$, and eigenvalue splittings and their context derivatives shrink like $e^{-\lambda t}$ with $\lambda=-f_p(p_*,\alpha_0)$. For $\alpha_0=\kappa=1$, $p_*=\sqrt2$ and the radial part of the tensor decays at rate $2\lambda=12-4\sqrt2\approx6.34$, as does the angular part when $\eta_0=\alpha_0$ (otherwise the angular part decays at twice the larger of the belief and model rates); by the root join it has lost more than three orders of magnitude. In the isotropic limit every rotation stabilizes the laws and the connection is undetermined. Keeping the geometry alive needs anisotropic local data in both sectors, context dependent in the belief sector, for example a belief evaluator $W_2(\alpha_0,\Gamma(a,\chi))$ with $\Gamma$ the prepared section and an anisotropic model reference; the precisions then relax to anisotropic fixed points and the curvature tends to a nonzero constant. The curvature is then a property of those declared data rather than of the interaction.

Context-constant individual shapes are substantive for the tensor, though not for the connection: an opposite contextual variation of two model shapes has positive fine Fisher length but zero derivative of their sum law, and any contextual variation of individual shapes, even a common one, destroys the claimed tensor equality, while the minimizing connection is still inherited. Common precision, a constant model-independent evaluator, global moment channels, product recognition, and fresh extensive responses likewise delimit this construction. The threshold is constitutive, and selection does not alter the underlying fine flow. \Cref{cf:open-formation} removes the common-parameter preparation through a different selective architecture with native preferential routing and retained contextual profiles. Grouping that depends on the interaction, a derived interface or membership criterion, and unrestricted child probes remain separate problems. Neither a thermodynamic limit nor spacetime follows from these finite constructions.
